\section{Settings}\label{sec:settings}

\subsection{Two settings objects}

Presentation --- label style, colourmap, time-axis convention, figure sizes ---
belongs to \code{pymorgan.Settings} and is read through
\code{pm.get\_settings()} after \code{pm.apply\_style()}. PyRATE-TA's own
\code{Settings} holds \emph{analysis-only} fields, listed in
Table~\ref{tab:settings}. There is deliberately no second aesthetics system.

\begin{table}[h]
  \centering
  \begin{tabular}{llll}
    \toprule
    Field & Default & Group & Meaning \\
    \midrule
    \code{n\_components}        & 2            & fit    & Components proposed for a new fit \\
    \code{model\_type}          & Sequential   & fit    & Default model family \\
    \code{irf\_mode}            & Gaussian     & fit    & \code{None} / \code{Gaussian} / \code{Clip} \\
    \code{irf\_fwhm}            & \code{None}  & fit    & Fixed IRF width; \code{None} fits it \\
    \code{t\_min}               & \code{None}  & fit    & Clip delays below this (mode \code{Clip}) \\
    \code{ftol}                 & $10^{-8}$    & solver & Termination tolerance on the cost \\
    \code{xtol}                 & $10^{-8}$    & solver & Termination tolerance on the parameters \\
    \code{max\_iterations}      & 2000         & solver & Cap on iterations; exceeding it is a failure \\
    \code{variable\_projection} & \code{True}  & solver & Solve amplitudes in closed form \\
    \code{report\_uncertainties}& \code{True}  & solver & Report linearised 1-$\sigma$ errors \\
    \code{use\_noise\_weights}  & \code{False} & solver & Weight residuals by the per-point noise \\
    \code{noise\_floor\_fraction} & $10^{-3}$  & solver & Clip $\sigma$ below this fraction of the median \\
    \code{degeneracy\_warn\_ratio} & 1.2       & solver & Warn when two lifetimes come this close \\
    \midrule
    \code{default\_datadir}     & \code{None}  & paths  & Directory offered by the file dialog \\
    \bottomrule
  \end{tabular}
  \caption{The analysis settings that appear in the settings dialog. LDA
    parameters are not listed here; they are set via the GUI LDA tab or
    directly in \code{settings.toml} for headless use
    (Table~\ref{tab:lda-settings}).}
  \label{tab:settings}
\end{table}

\begin{table}[h]
  \centering
  \begin{tabular}{llll}
    \toprule
    Field & Default & Meaning \\
    \midrule
    \code{lda\_n\_lifetimes}      & 100          & Grid points in the lifetime axis \\
    \code{lda\_tau\_min}          & \code{None}  & Grid lower bound; \code{None} = auto \\
    \code{lda\_tau\_max}          & \code{None}  & Grid upper bound; \code{None} = auto \\
    \code{lda\_regularisation}   & L-curve      & Alpha selection: \code{L-curve} / \code{GCV} / \code{Fixed} \\
    \code{lda\_alpha}            & $10^{-3}$    & Strength used when \code{lda\_regularisation} is \code{Fixed} \\
    \code{lda\_non\_negative}    & \code{False} & Constrain distribution to $\ge 0$ \\
    \bottomrule
  \end{tabular}
  \caption{LDA default settings, stored in the \code{[lda]} section of
    \code{settings.toml}. These do \emph{not} appear in the settings dialog;
    the GUI LDA tab reads them at start-up to pre-populate its widgets, which
    are then the single control surface for an interactive session.}
  \label{tab:lda-settings}
\end{table}

Turning \code{variable\_projection} off is almost always the wrong choice
(section~\ref{sec:varpro}); it is exposed for diagnostics only.

\subsection{Design}

\code{Settings} is a dataclass and \textbf{the dataclass field is the single
source of truth}. \code{\_\_post\_init\_\_} (which coerces enums),
\code{to\_dict}, \code{from\_dict} and \code{update\_settings} all derive from
\code{dataclasses.fields()}, so adding a setting needs only the field. Enums are
\code{enum.StrEnum}: \code{str(member)} is the value, so they serialise and
compare as plain strings while staying type-checked.

\code{Settings.field\_specs()} drives the GUI settings panel: \code{label},
\code{kind} (\code{choice} / \code{float} / \code{int} / \code{bool} /
\code{text}), optional \code{choices} / \code{min} / \code{max} /
\code{step} / \code{tooltip}, and \code{tab}. The method is \emph{derived}
rather than hand-written: every dataclass field gets an entry automatically,
with the widget kind read from its annotation, the page from
\code{field\_sections()} and the tooltip from \code{field\_comments()}.
\emph{A setting without an entry there does not appear in the panel} ---
which is the intended way to hide an internal knob, and the usual explanation
for a new setting not showing up. A small override table refines labels and
ranges, and can hide a deprecated field.

\subsection{Sections of the file}

\code{settings.toml} is grouped by concern, with a comment above every key:
\code{[fit]} (defaults for a new fit), \code{[solver]} (optimiser limits,
weighting, reporting), \code{[irf]}, \code{[lda]} (LDA defaults for headless
and scripted use), \code{[gui]} and \code{[paths]}. The mapping from field to
section is \code{Settings.field\_sections()}, and it is exhaustive: a field
with no section raises on save rather than being written where the loader would
not find it.

\paragraph{The \code{[lda]} block and the two control surfaces.}
LDA has two ways to set its parameters:
\begin{itemize}
  \item \textbf{Interactive}: the GUI LDA tab --- dropdowns, spinboxes and
    checkboxes. These are the canonical control surface during a session.
  \item \textbf{Headless / scripted}: \code{settings.toml} under \code{[lda]}.
    Keyword arguments to \code{solve\_lda} override everything for a single
    call.
\end{itemize}
\noindent
The \code{[lda]} block does \emph{not} appear in the settings dialog, because
that would create two places to set the same parameters. Instead, the GUI reads
\code{lda\_*} fields from the active \code{Settings} once at start-up and
loads them into the tab's widgets (\code{LDATabMixin.\_init\_lda\_tab\_from\_settings}),
so editing \code{settings.toml} changes the defaults the GUI sees the next time
it opens.

\noindent
The same precedence applies for scripted use:

\begin{lstlisting}
# defaults from settings.toml
result = pr.fit.lda.solve_lda(data)

# or override explicitly --- keyword args always win
result = pr.fit.lda.solve_lda(
    data,
    n_taus=200,
    alpha="auto",
    alpha_method="gcv",
    non_negative=True,
)
\end{lstlisting}

A flat file --- the older layout --- is still read, so an existing
configuration is never orphaned.

\subsection{The configuration file}

Settings are read from and written to \code{settings.toml} with \code{tomlkit},
so user comments survive a save. The file is looked up in the working
directory first, then at the project root.

\begin{lstlisting}
import pyrate_ta as pr

pr.load_settings("settings.toml")
pr.update_settings(n_components=3, irf_mode="Clip", t_min=0.3)
pr.save_settings("settings.toml")
\end{lstlisting}

\code{update\_settings} rejects unknown names rather than storing them, so a
typo fails immediately instead of silently doing nothing.

\subsection{The settings editor}

\code{uv run pyrate-ta-settings} opens a standalone editor whose pages and fields
are generated from \code{field\_specs()}; \emph{Save permanently} writes them
back to \code{settings.toml}. A value the settings layer rejects is logged and
reverted in the widget rather than silently dropped. \code{settings.py} itself
never imports Qt.

The same editor is reachable from the main window under
\emph{View $\rightarrow$ Settings...}, in a dialog with two halves:
\emph{Analysis (PyRATE-TA)} and \emph{Aesthetics (PyMORGAN)}. Both panels edit
their own package's settings live --- the embedded plots re-render as values
change --- and \emph{Save permanently} writes each half back to \emph{its own}
\code{settings.toml}. PyMORGAN's file is located from the installed package
rather than from the working directory: when PyRATE-TA is the host, the
\code{settings.toml} in the working directory is PyRATE-TA's, and writing
PyMORGAN's fields into it would overwrite the wrong file.

\subsection{Both files are loaded}

The GUI loads PyRATE-TA's \code{settings.toml} \emph{and} PyMORGAN's. Without the
second, PyMORGAN's settings stayed at their dataclass defaults, so a profile or
a \code{font\_scale} chosen in PyMORGAN was ignored here and the same data was
drawn with a different font stack --- which is what makes a PyRATE-TA figure look
unlike a PyMORGAN one while both claim the same style. PyMORGAN's file is
located from the installed package, never from the working directory, for the
reason given above.

Fonts themselves are PyMORGAN's (\code{Helvetica, TeX Gyre Heros, Arial}, from
its \code{.mplstyle}) and PyRATE-TA applies them unchanged, but the
\emph{installation} is per environment: a PyRATE-TA virtual environment that has
never run \code{pymorgan-install-fonts} falls back to DejaVu Sans. The GUI
therefore checks at start-up and says so, with the command to run.

\subsection{How results are quoted}

Two settings in \code{[plots]} govern every lifetime PyRATE-TA prints:

\begin{description}
  \item[\code{show\_uncertainties}] Whether the $\pm$ appears at all, in the
    plot legends and the logged fit summary.
  \item[\code{round\_uncertainties}] Round a value \emph{together with} its
    uncertainty: the uncertainty to one significant figure and the value to
    \emph{that same decimal place}, so $325.3 \pm 2.5$ is quoted as
    $325 \pm 3$, $297.5 \pm 1.38$ as $298 \pm 1$ and $8.036 \pm 0.0234$ as
    $8.04 \pm 0.02$. This keeps a fitted constant from being presented with
    more precision than its own uncertainty supports. Halves round \emph{up}
    ($2.5 \rightarrow 3$): an uncertainty should never be rounded down.
\end{description}

Both are applied at the point of display only; the result object keeps the full
precision, so a saved session and the copied lifetimes are unaffected.
